% =============================================================
% latex/zeichenfeld.tex
%
% \zeichenfeld[<ean>]{<feld-nr>}
%
%   Rendert ein Werkstatt-Feld zum Selber-Zeichnen eines EAN-13-
%   Strichcodes. Bestandteile:
%     - Kopfzeile mit Feld-Nummer und EAN-Klartext (leer oder
%       eingetragen)
%     - Hilfstabelle: pro Position 1..13 die Spalten Ziffer, Set
%       (L/G/R) und 7 Bits, entweder leer oder vorausgefüllt
%     - Zeichenfeld in echter EAN-Größe (1 Modul = 1mm), mit
%       feinem Modul-Hilfsraster und Markern für die Bauteil-
%       grenzen (Hellzonen, Guards, Datenstellen).
%     - Klarschrift-Zeile unter dem Code (Kästchen oder gefüllt)
%
%   Wenn das optionale [<ean>] mitgegeben wird, wird der Bogen
%   als Lösung gerendert (Tabelle ausgefüllt, Strichcode in den
%   Modulen schwarz gemalt). Ohne Argument bleibt alles leer und
%   muss mit Bleistift/Fineliner ausgefüllt werden.
%
% Voraussetzungen: TikZ, xstring, ean13-encoding (alle in
% praeambel.tex geladen).
% =============================================================

\makeatletter

% --- Hilfstabelle: 13 Spalten, 4 Zeilen, breit für Bits ---
% Spaltenbreite 11mm, y-Skala großzügig damit Hand-Schrift reinpasst.
% Bits sind in der ausgefüllten Variante 45° gedreht (sonst kollidieren
% die 7 Zeichen mit der Nachbarspalte).
\NewDocumentCommand{\@zfTabelle}{m}{%
  \begin{tikzpicture}[x=11mm, y=6mm, baseline=0]
    % Zeilen-Beschriftung links
    \node[anchor=east, font=\small\bfseries] at (0.4,4)   {Position};
    \node[anchor=east, font=\small\bfseries] at (0.4,3)   {Ziffer};
    \node[anchor=east, font=\small\bfseries] at (0.4,2)   {Set};
    \node[anchor=east, font=\small\bfseries] at (0.4,0.6) {7 Bits};
    %
    % Position 1 ist die Erstziffer — sie hat eine Ziffer, aber kein
    % Set und keine 7 Bits (sie wird nicht direkt codiert, sondern
    % steckt im Paritätsmuster der Folge-Ziffern). Wir zeichnen nur
    % das Ziffer-Feld und markieren Set + Bits mit „—".
    \node[font=\small] at (1,4) {1};
    \IfStrEq{#1}{}{%
      \draw[gray] (1-0.4, 2.55) rectangle (1+0.4, 3.45);% Ziffer
    }{%
      \StrChar{#1}{1}[\@zfZ]%
      \node[font=\small] at (1,3) {\@zfZ};
    }
    \node[font=\small, text=gray] at (1,2)   {---};
    \node[font=\small, text=gray] at (1,0.6) {---};
    %
    % Positionen 2..13: voll bzw. leere Kästchen.
    \foreach \p in {2,...,13}{%
      \node[font=\small] at (\p,4) {\p};
      \IfStrEq{#1}{}{%
        \draw[gray] (\p-0.4, 2.55) rectangle (\p+0.4, 3.45);% Ziffer
        \draw[gray] (\p-0.35, 1.55) rectangle (\p+0.35, 2.45);% Set
        \draw[gray] (\p-0.45, -0.2) -- (\p+0.45, 1.4);% Bits-Schreiblinie
      }{%
        \StrChar{#1}{\p}[\@zfZ]%
        \node[font=\small] at (\p,3) {\@zfZ};
        \node[font=\small\ttfamily] at (\p,2) {\eantableset{#1}{\p}};
        \node[font=\small\ttfamily, rotate=45] at (\p,0.6) {\eantablebits{#1}{\p}};
      }
    }%
    \draw[gray, line width=0.3pt] (0.5, 3.6) -- (13.5, 3.6);
  \end{tikzpicture}%
}

% --- Strichcode-Zeichenbereich mit Raster ---
\NewDocumentCommand{\@zfCode}{m}{%
  \begin{tikzpicture}[x=1mm, y=1mm, baseline=0]
    % Hellzonen-Rahmen (gestrichelt, sehr dezent)
    \draw[gray!50, dashed, line width=0.2pt] (-11,-3) rectangle (102,33);
    %
    % Sehr feines Modul-Hilfsraster: vertikale Linien an Modul-Grenzen
    \foreach \x in {0,1,...,95}{%
      \draw[gray!25, line width=0.05pt] (\x,-3) -- (\x,33);%
    }
    % Untere und obere Begrenzung der Daten-Strichhöhe (y=0..33)
    \draw[gray!40, line width=0.15pt] (0,0) -- (95,0);
    \draw[gray!40, line width=0.15pt] (0,33) -- (95,33);
    %
    % Marker für Bauteil-Grenzen (oben drüber, zwei Reihen damit nichts kollidiert)
    % Reihe oben: Bezeichner der Sechsergruppen
    \node[anchor=south, font=\scriptsize\sffamily] at (24,38)     {Ziffern 2--7 (links)};
    \node[anchor=south, font=\scriptsize\sffamily] at (71,38)     {Ziffern 8--13 (rechts)};
    % Reihe unten: Guards
    \node[anchor=south, font=\scriptsize\sffamily] at (1.5,33.5)  {Start};
    \node[anchor=south, font=\scriptsize\sffamily] at (47.5,33.5) {Mitte};
    \node[anchor=south, font=\scriptsize\sffamily] at (93.5,33.5) {Ende};
    %
    % Vertikale Strichlinien zur Bauteil-Trennung (etwas kräftiger,
    % über den Strichbereich nach oben/unten hinaus)
    \draw[gray!60, line width=0.4pt] (0,-5) -- (0,36);   % Start-Guard begin
    \draw[gray!60, line width=0.4pt] (3,-5) -- (3,36);   % Start-Guard end
    \draw[gray!60, line width=0.4pt] (45,-5) -- (45,36); % Mittel-Guard begin
    \draw[gray!60, line width=0.4pt] (50,-5) -- (50,36); % Mittel-Guard end
    \draw[gray!60, line width=0.4pt] (92,-5) -- (92,36); % End-Guard begin
    \draw[gray!60, line width=0.4pt] (95,-5) -- (95,36); % End-Guard end
    %
    % Ziffer-Trennungen — alle 7 Module innerhalb der Datenbereiche.
    % Auch nach oben/unten verlängert, damit das Zähl-Raster für
    % die Sechsergruppen optisch sofort erkennbar ist.
    \foreach \x in {10,17,24,31,38, 57,64,71,78,85}{%
      \draw[gray!50, line width=0.25pt] (\x,-5) -- (\x,36);%
    }
    %
    % Wenn EAN mitgegeben: Module schwarz füllen
    \IfStrEq{#1}{}{}{%
      \eanbits{#1}{\@zfBits}%
      \foreach \bit [count=\j from 0] in \@zfBits {%
        \ifnum\bit=1
          \def\@zfBo{0}%
          \ifnum\j<3\relax\def\@zfBo{-3}\fi
          \ifnum\j>91\relax\def\@zfBo{-3}\fi
          \ifnum\j>44\relax\ifnum\j<50\relax\def\@zfBo{-3}\fi\fi
          \fill[black] (\j,\@zfBo) rectangle ({\j+1},33);%
        \fi
      }%
    }%
    %
    % Klarschrift-Bereich unten: Kästchen oder gefüllt
    \IfStrEq{#1}{}{%
      % leere Kästchen für 1 + 6 + 6 Ziffern
      \draw[gray] (-9,-8) rectangle (-2,-3);
      \foreach \i in {0,...,5}{\draw[gray] (3+\i*7, -8) rectangle (3+\i*7+7, -3);}%
      \foreach \i in {0,...,5}{\draw[gray] (50+\i*7, -8) rectangle (50+\i*7+7, -3);}%
    }{%
      % Voll: Ziffern an denselben Positionen wie die Kästchen oben,
      % damit der Mustercode optisch zur leeren Vorlage passt.
      \StrChar{#1}{1}[\@zfF]%
      \node[anchor=center, font=\sffamily\large] at (-5.5,-5.5) {\@zfF};
      \foreach \i in {0,...,5}{%
        \StrChar{#1}{\the\numexpr\i+2\relax}[\@zfDi]%
        \node[anchor=center, font=\sffamily\large] at ({3+\i*7+3.5},-5.5) {\@zfDi};%
      }%
      \foreach \i in {0,...,5}{%
        \StrChar{#1}{\the\numexpr\i+8\relax}[\@zfDi]%
        \node[anchor=center, font=\sffamily\large] at ({50+\i*7+3.5},-5.5) {\@zfDi};%
      }%
    }%
  \end{tikzpicture}%
}

% --- Komplettes Feld ---
\NewDocumentCommand{\zeichenfeld}{O{} m}{%
  \par\noindent
  \begin{minipage}{\linewidth}
    \rule{\linewidth}{0.4pt}\par\smallskip
    \textbf{\large Feld #2}\hfill
    \textbf{EAN:}~\IfStrEq{#1}{}{%
      \framebox[7em]{\strut}%
    }{%
      \texttt{#1}%
    }\par
    \medskip
    \centerline{\@zfTabelle{#1}}
    \medskip
    \centerline{\@zfCode{#1}}
  \end{minipage}\par
  \vspace{1em}
}

% --- Cheat-Sheet: Codetabellen + Paritätsmuster, kompakt ---
% Soll am Fuß der Werkstatt-Seite stehen, damit Greta beim Ausfüllen
% nichts nachschlagen muss.
\NewDocumentCommand{\zeichenfeldcheatsheet}{}{%
  \par\smallskip
  \noindent\textbf{\small Spickzettel — Codetabellen und Paritätsmuster}\par
  \smallskip
  \noindent
  \begin{minipage}[t]{0.62\linewidth}
    \footnotesize
    \setlength{\tabcolsep}{4pt}
    \begin{tabular}{c|c|c|c}
      \toprule
      \textbf{Ziffer} & \textbf{L-Code} & \textbf{G-Code} & \textbf{R-Code} \\
      \midrule
      0 & \texttt{0001101} & \texttt{0100111} & \texttt{1110010} \\
      1 & \texttt{0011001} & \texttt{0110011} & \texttt{1100110} \\
      2 & \texttt{0010011} & \texttt{0011011} & \texttt{1101100} \\
      3 & \texttt{0111101} & \texttt{0100001} & \texttt{1000010} \\
      4 & \texttt{0100011} & \texttt{0011101} & \texttt{1011100} \\
      5 & \texttt{0110001} & \texttt{0111001} & \texttt{1001110} \\
      6 & \texttt{0101111} & \texttt{0000101} & \texttt{1010000} \\
      7 & \texttt{0111011} & \texttt{0010001} & \texttt{1000100} \\
      8 & \texttt{0110111} & \texttt{0001001} & \texttt{1001000} \\
      9 & \texttt{0001011} & \texttt{0010111} & \texttt{1110100} \\
      \bottomrule
    \end{tabular}
  \end{minipage}\hfill
  \begin{minipage}[t]{0.35\linewidth}
    \footnotesize
    \setlength{\tabcolsep}{4pt}
    \begin{tabular}{c|l}
      \toprule
      \textbf{Erstziffer} & \textbf{Pattern} \\
      \midrule
      0 & \texttt{LLLLLL} \\
      1 & \texttt{LLGLGG} \\
      2 & \texttt{LLGGLG} \\
      3 & \texttt{LLGGGL} \\
      4 & \texttt{LGLLGG} \\
      5 & \texttt{LGGLLG} \\
      6 & \texttt{LGGGLL} \\
      7 & \texttt{LGLGLG} \\
      8 & \texttt{LGLGGL} \\
      9 & \texttt{LGGLGL} \\
      \bottomrule
    \end{tabular}\par\smallskip
    \scriptsize\itshape
    Position 8--13 (rechte Hälfte) immer R-Code.
  \end{minipage}%
  \par\medskip
}

\makeatother
